\documentclass[11pt]{article}

\usepackage[T1]{fontenc}
\usepackage[utf8]{inputenc}
\usepackage[margin=1in]{geometry}
\usepackage{microtype}
\usepackage{amsmath,amssymb,amsthm}
\usepackage{graphicx}
\usepackage{booktabs}
\usepackage{xcolor}
\usepackage{listings}
\usepackage[hidelinks]{hyperref}
\usepackage[nameinlink,noabbrev]{cleveref}

\newtheorem{theorem}{Theorem}
\newtheorem{definition}{Definition}

\lstset{
  basicstyle=\ttfamily\small,
  frame=single,
  columns=fullflexible,
  keepspaces=true
}

\title{Auditable Editing for Structured Research Papers}
\author{
  Mira Chen \\ Northstar Research \\ \texttt{mira@example.edu}
  \and
  Tomas Rivera \\ Northstar Research \\ \texttt{tomas@example.edu}
}
\date{}

\begin{document}

\maketitle

\begin{abstract}
Research editors must balance structured interaction with exact source
preservation. We describe a synthetic benchmark for small, revision-checked
LaTeX edits and report illustrative results on a deliberately modest corpus.
The example demonstrates Setwright's supported paper structures; it does not
make an empirical product claim.
\end{abstract}

\section{Introduction}
\label{sec:introduction}

LaTeX gives researchers durable, inspectable source, while structured editors
can reduce the cost of common operations. Existing scientific workflows also
depend on citations, stable cross-references, and mathematical notation
\cite{lamport1994latex,vaswani2017attention}. A dependable bridge between those
interfaces cannot silently rewrite syntax it does not understand.

We study a narrow question: can a structured edit be represented as a minimal
byte patch and rejected when its source precondition is stale? This sample uses
a static include graph so every visual operation remains inside one file
boundary. A source span is measured in UTF-8 bytes, not displayed characters.\footnote{This distinction matters for non-ASCII source text.}

Our illustrative pipeline is shown in \cref{fig:pipeline}; the validation rule
is formalized in \cref{sec:method}.

\begin{figure}[t]
  \centering
  \fbox{\parbox[c][2.6cm][c]{0.78\linewidth}{\centering
    Intent $\rightarrow$ minimal patch $\rightarrow$ candidate parse
    $\rightarrow$ atomic commit}}
  \caption{A source-authoritative editing pipeline. The diagram is rendered
  with ordinary LaTeX so the sample has no binary asset dependency.}
  \label{fig:pipeline}
\end{figure}

\section{Method}
\label{sec:method}

\begin{definition}[Validated edit]
A validated edit is a tuple
$e=(f,b_0,b_1,r,h)$ containing a file identifier, half-open byte interval,
replacement text, and the expected hash of the replaced bytes.
\end{definition}

Let $S_f$ be the authoritative bytes for file $f$. An edit is eligible exactly
when its interval is valid and its precondition matches:
\begin{equation}
  \label{eq:eligibility}
  0 \le b_0 \le b_1 \le \lvert S_f \rvert
  \quad\land\quad
  H(S_f[b_0:b_1]) = h.
\end{equation}

\begin{theorem}[Stale-edit exclusion]
If the bytes inside an edit's interval change after the edit is proposed, a
collision-resistant expected-slice hash causes validation to reject the edit,
except with negligible collision probability.
\end{theorem}

\begin{proof}
Changing the interval changes the hash input. Acceptance would therefore
require a collision or second preimage for the selected digest.
\end{proof}

The benchmark applies an edit, reparses the candidate, checks complete source
coverage, and then applies the inverse edit. A compact representation is:

\begin{lstlisting}[language={},caption={Illustrative edit loop.},label={lst:loop}]
for edit in edits:
    candidate = validate_and_apply(source, edit)
    assert complete_coverage(candidate)
    source = apply_inverse(candidate, edit)
\end{lstlisting}

\section{Illustrative results}
\label{sec:results}

\Cref{tab:results} contains synthetic values chosen only to exercise table
editing and references. They are not measurements of the Setwright scaffold.

\begin{table}[t]
  \centering
  \caption{Synthetic edit-validation outcomes.}
  \label{tab:results}
  \begin{tabular}{lrr}
    \toprule
    Fixture class & Accepted & Correctly rejected \\
    \midrule
    Plain paragraphs & 120 & 0 \\
    Equations & 84 & 6 \\
    Raw environments & 0 & 32 \\
    Stale source spans & 0 & 40 \\
    \bottomrule
  \end{tabular}
\end{table}

The intended interpretation follows \cref{eq:eligibility}: ordinary edits with
matching preconditions proceed, while stale or intentionally unsupported edits
remain explicit. A real evaluation would report corpus licenses, hardware,
software versions, confidence intervals, and failure cases.


\section{Limitations and conclusion}
\label{sec:conclusion}

The benchmark omits arbitrary macro expansion and stateful TeX. Those regions
must remain exact raw source rather than being normalized into a visual model.
Future work should evaluate larger, permissively licensed paper corpora and all
supported operating-system sandboxes.

\bibliographystyle{plain}
\bibliography{references}

\end{document}
